% SPDX-License-Identifier: GPL-3.0-or-later
% English -- what the reading editions' preamble leaves to the language.
% Input by lib/frank-preamble.tex as \input{\FrankLib/lang/\BookLang.tex};
% docs/languages.md section 6 lists the hooks every file here must define and
% section 12 the ones it may.
%
% The Latin-script model, as it.tex: no font of its own, no marks to strip,
% two passes.  What is English's own is that it is also the language the
% GLOSSES are written in.  Every other file here describes a language the
% core preamble treats as foreign; this one describes the language the core
% preamble is already in.  So \FrankLangName below is the same name the core
% hands \begin{otherlanguage}{english} for every gloss block, and an English
% edition is the first where the chunk column and the gloss column are set in
% one language.  Nothing in the core asks whether they differ, which is why
% this costs no code -- but it is why \pw is not used on this page (below),
% and why the middle line of the gloss carries more weight here than
% anywhere else.

% ---- the babel language and the switches ----------------------------------
% \FrankLangName is what \foreignlanguage and \begin{otherlanguage} are given
% everywhere in the core preamble: babel must know this name, or every passage
% comes out in the roman with a "Unknown language" error.  Here it is the name
% babel was loaded with in the first place.
\newcommand{\FrankLangName}{english}
\FrankRTLfalse           % chunk left, gloss right; \pw needs no bidi isolate
\FrankHasBarefalse       % nothing is written in, so nothing can come off again
\FrankHasAltfalse        % no alternate face
\FrankHasReadingfalse
\FrankHasVertfalse

% ---- fonts ------------------------------------------------------------------
% The registry's tex.main is null for English: the text is set in the main
% roman the core preamble already loaded (TeX Gyre Pagella), which is what a
% Latin script wants and here is not even a choice -- the glosses are English
% too, so a face of its own would set the same language two ways on one page.
%
% import=en is not redundant, though the core already loaded babel with the
% `english' option.  Measured with \showhyphens{establishment representative}
% under lualatex, with and without this line: both give es-tab-lish-ment
% rep-re-sen-ta-tive, so the break points are untouched -- what the line adds
% is one Info in the log, "Importing data for english from babel-en.ini", and
% with it the locale data (the quotation marks, the case pairs) comes from the
% same ini file every other language's does instead of from english.ldf alone.
% Hyphenation is why that had to be checked and not assumed: hyphen-english is
% in every TeX Live, so unlike the other Latin four English cannot lose its
% patterns -- and a second \babelprovide silently replacing them would be the
% one way it could.
\babelprovide[import=en]{english}
% The pronunciation line is IPA (docs/lang/en.md), so this language leans on
% the core's DejaVu Serif fallback harder than any other: asked of luaotfload
% on this machine, TeX Gyre Pagella has theta, eth, eng, ae and schwa and has
% NO glyph for esh, ezh, open-e, script-a, small-capital-I, upsilon, turned-v,
% open-o, the stress mark or the length mark -- ten of the sixteen letters a
% transcription needs.  DejaVu Serif has all sixteen.  A machine without it
% still builds, and every one of those ten is then a "Missing character" line
% in the .log and a hole in the page, which is the same bargain Italian and
% French already take for a smaller number of letters.

% ---- the vocabulary line -----------------------------------------------------
% \vb{plain form}{sound}{past}{sound}{past participle}{sound}{meaning}: the
% three principal parts, which is what an English dictionary prints for an
% irregular verb and the only thing about a verb that cannot be worked out
% (go, went, gone; regular verbs make both with -ed, and are given all three
% anyway so the slots always mean the same thing).  The present is not among
% them: English inflects it in one person only -- except in four verbs, and
% those four get it in ONE parenthesis after the meaning, a closed list:
%   be    (pres. \pw{am}, \pw{is}, \pw{are})     have  (pres. \pw{has})
%   do    (pres. \pw{does} \textit{dʌz})        say   (pres. \pw{says} \textit{sɛz})
% -- does and says because their vowel is not the one the spelling promises.
% Nothing else goes in a parenthesis for English.  A MODAL (can could will
% would shall should may might must ought) gets no \vb at all: it has no
% principal parts to give, "p.p. could" would be a form that does not exist,
% so it is a \dw and its past, where it has one, is named in the meaning.
% A phrasal verb is one headword in all three slots (give up, gave up,
% given up) however far the sentence has carried the particle.
% The reader takes the same pair from the registry (lib/languages.json,
% "vb_labels": ["past", "p.p."], the three forms in words as "vb_forms"); the
% two must agree, or the PDF and the reader of one book label the same form
% differently.  lib/newlang.py --check compares them.
\newcommand{\FrankVbPres}{past}
\newcommand{\FrankVbPast}{p.p.}

% ---- the two Lua filters ---------------------------------------------------
% frank_strip(s): the bare form of a chunk -- what pass 3 shows.  Written from
% the registry's `strip`; the identity when there is nothing to take off.
% frank_latin(s): the label's digits as ASCII, because a PDF destination name
% cannot carry the language's own figures.  Written from `digits`.
% Both are the identity for English: nothing is added to the page for pass 1
% to carry, and the figures are already 0-9.  They still have to exist -- the
% core calls them for every language.  No #, % or literal backslash in a
% \directlua body: TeX reads it first (NOTES section 9 of the Persian book).
\directlua{
  function frank_strip(s) return s end
  function frank_latin(s) return s end
}

% ---- the "How to read this" page -----------------------------------------------
% Printed by lib/frank-frontmatter.tex, before the first chapter: the only
% place the reader is told what the passes are for, and the only place the
% alphabet of the middle line is explained.
% English words are set in \textit and not in \pw, as in it.tex and de.tex.
% \pw shrinks a word to 8.6pt so a foreign script can be told from the English
% around it; here the word and the prose are the same language in the same
% face, and shrinking it would say nothing.  The sounds are in \textit too --
% that is what the gloss line itself does with them -- and they need no other
% mark to be told from the words, because the letters they are written in are
% not English letters.  That is the whole argument for the alphabet, and the
% page makes it by using it.
\newcommand{\FrankHowTo}{%
Every passage comes twice.

\smallskip
\textbf{First} the whole sentence, and nothing else. Try to read it, aloud if
you can. Get as far as you can and do not worry about what you cannot get —
the point is the attempt, made before any help arrives.

\smallskip
\textbf{Second} the same sentence in small chunks, English on the left, and on
the right three lines: how it is said, then how to look each word up, then
what it means. Now you find out how much you had right. There is no third
pass: nothing was written into the English for the first one, so there is
nothing to take away.

\smallskip
The middle line is here for every chunk, which it is in no other edition of
this series: \textit{through}, \textit{though}, \textit{thought} and
\textit{tough} share five letters and no vowel. It is written in the
dictionaries' alphabet, because respelled English would be read as English:
\textit{ˈbʌtər} cannot be mistaken for a spelling of \textit{butter}.

\smallskip
\textit{ˈ} stands before the stressed syllable, and every word longer than one
has it: stress is what English does not write and the thing most often got
wrong (\textit{ˈrɛkərd} the thing, \textit{rəˈkɔrd} the act).
\textit{ə} is the colourless vowel of every unstressed syllable — the
\textit{a} of \textit{about}, the \textit{o} of \textit{lemon}. The vowels that
keep their colour are \textit{æ} \textit{cat}, \textit{ɛ} \textit{bed},
\textit{ɪ} \textit{sit}, \textit{ɑ} \textit{hot}, \textit{ʌ} \textit{cup},
\textit{ʊ} \textit{book}, \textit{ɔ} \textit{caught}, \textit{i}
\textit{see}, \textit{u} \textit{too}; \textit{θ} \textit{ð} \textit{ʃ}
\textit{ʒ} \textit{ŋ} the sounds of \textit{thin}, \textit{this},
\textit{sh}, \textit{measure}, \textit{sing}; \textit{j} the \textit{y} of
\textit{yes}, \textit{dʒ} the \textit{j} of \textit{jam}. The rest are as you
expect, and every written \textit{r} is sounded: the accent is General
American.

\smallskip
Verbs are given as plain form, past, past participle — \textit{go},
\textit{went}, \textit{gone}. One with a particle is given whole,
\textit{give up}, however far the sentence has carried it: the meaning is the
pair's, never the sum. Four verbs have a present you could not build, and it
stands in brackets after the meaning: \textit{am}, \textit{is}, \textit{are}
for \textit{be}; \textit{has}; \textit{does} \textit{dʌz}; \textit{says}
\textit{sɛz}. \textit{Can}, \textit{must}, \textit{will} and the other
modals have none of these forms, and are given as words like any other.
}
